\newpage

\section{Obstacle Evasion Course}

The "Obstacle Evasion Course" tests the vehicle's advanced driving capabilities and decision-making in complex traffic scenarios.

% TODO: Removal of any kind of mention of rural road and suburban area throughout the document

% The “Obstacle Evasion Course” extends the track of the Free Drive event with additional elements which need to be considered during the driving task.
% Static and dynamic obstacles are added to the rural road scenario.
% Additionally, the vehicle has to perform automatic parking maneuvers by finding a suitable parking spot inside a parking lot.
% The track does additionally contain at least one suburban section at this point.
% All definitions concerning the course of the road maintain validity.

% Oncoming traffic is not to be expected, except when passing barred areas inside the suburban scenario.

\subsection{Scenario}

The track of the Obstacle Evasion Course can contain any of the elements described in Section \ref{track_layout} except for landmarks.

Static and dynamic obstacles are present on the track, intersections may specify turning directions and right of way regulations.
All traffic signs defined in Section \ref{traffic_signs} can be present and must be respected.

The vehicle is expected to stay in the right driving lane at all times.
Crossing the lane markings with two or more wheels will result in a penalty.

The vehicle will be evaluated for each element it passes according to the scoring guidelines in Section \ref{obstacle_scoring_guidelines}.
Additionally, parking maneuvers can be performed once per lap.

Elements will not be located on uphill and downhill grades. 

\subsection{Elements of the Obstacle Evasion Course}

The following sections describe the vehicle's expected behavior when encountering specific elements on the track.

\label{elements_obstacle_evasion}
\subsubsection{Static and Dynamic Obstacles}

Obstacles may force the vehicle to change lanes. 
Lane changes must be indicated using the turn indicators.
Passing maneuvers must be executed without touching an obstacle.
After the obstacle has been passed, the vehicle must return to the right driving lane within 2 meters.

Apart from static obstacles, at least one dynamic obstacle will be present on the track.
It may be passed, but not in intersections.
Passing maneuvers in intersections are penalized.

\subsubsection{Intersections}
\label{intersection_rural}

% Intersections of the rural road scenario must be crossed driving straight.
Intersections without a mandatory turning direction must be crossed driving straight.
A mandatory turning direction is indicated by the lane markings and a corresponding traffic sign (cf. Sections \ref{intersection} and \ref{traffic_signs}).

If a stop line is located in the driving lane, the vehicle must stop for at least 3 seconds.
The front of the vehicle may not cross the stop line, but must be positioned no further than 150 mm away from it.

At a give-way line, the vehicle must stop for at least 1 second.
The same regulations as for stop lines apply.

The right of way of a dynamic obstacle must be respected at an intersection, if the dynamic obstacle is located within the defined area (cf. Section \ref{fig_intersection_give_way}).
If the vehicle does not possess the right of way, it must wait until the dynamic obstacle has completely crossed the intersection.

\subsubsection{No-Passing Zones}
\label{event_no_passing_zones}

In a no-passing zone, the dynamic obstacle must be followed with a distance of at least 300 mm until the end of the zone.

\subsubsection{Two-lane Expressway}

Sections of the rural scenario, can be defined as an expressway.
Since the track is assumed to be a two-lane expressway, the vehicles must stay in the right lane all the time. \HighlightFix{Ref Issue}

\subsubsection{Barred Area}

A barred area may not be entered by the vehicle.
It must be passed by switching to the left lane if necessary.
The same regulations as for overtaking obstacles apply.

If a dynamic obstacle is located within 1000 mm of the beginning of the barred area, the vehicle has to wait for the obstacle to pass.
Switching lanes is only allowed with an empty left lane, oncoming traffic must have completely passed.
The desired passing maneuver has to be indicated while waiting by flashing the left turn indicators.
% Only one dynamic obstacle at a time can occur at a barred area.
If the vehicle is able to pass a barred area without leaving the own lane or driving over the markings, the vehicle may continue along the barred area even in case of oncoming traffic.

\subsubsection{Crosswalk}

One or more crosswalks may be present on the track with an unknown number of pedestrians.
If at least one pedestrian is present in the defined zones, the vehicle must stop in front of the crosswalk.
Stopping must be performed with the same regulations as at intersections.
Pedestrians start crossing only after the vehicle has stopped.
If all relevant pedestrians have crossed in front of the vehicle, the vehicle may continue.
Driving on before all pedestrians start to cross and have cleared the crosswalk will be penalized.
If no pedestrian is present, the vehicle may slow down, but is not allowed to come to a complete stop.

\subsubsection{Parking}

Parking can be performed once per lap during the Obstacle Evasion Course.

The vehicle can perform a parking attempt by finding a parking spot within the parking areas and maneuvering into it, without touching surrounding obstacles.
Attempting more than one parking attempt per lap will be penalized.
Parking spots are indicated by the corresponding traffic sign.

The start of a parking maneuver must be signaled using the turn indicators.
The turn indicators must point in the direction of the parking spot (i.e. blink left if the parking spot is to the left of the vehicle).

A complete parking maneuver requires the vehicle to come to a full stop being located in a valid spot with all of its wheels and flash all turn indicators at least one time.
While maneuvering out of the parking spot, the vehicle may cross the left lane, but has to continue driving in the right lane after the parking maneuver.

The correct position of the vehicle will be checked from both sides of the track with the first flashing of the indicators.
Penalties may apply until the vehicle is driving with a positive speed along the route while touching the right lane with at least 3 wheels.

Leaving the outer boundaries of the parking area is penalized the same as leaving the right lane while driving.
The speed in the parking lot is limited to the speed on the adjacent road. Collisions with obstacles during the parking maneuver will be penalized.

Vehicles may move forward or backward into the parking space.
The left lane of the track may only be crossed during the actual parking maneuver.
When searching for a parking spot, the vehicle must continue to use the right lane.

\subsection{RC-Mode}

Skipping challenges of the obstacle course (by not driving in the right lane) will be punished with the \HighlightNew{neutral} penalty designated for the respective element.

\HighlightNew{Each obstacle course element may be evaluated only once per round.
Elements may not be retried using RC-mode.
If incorrect behavior with respect to an element occurs before the activation of RC-mode, the element will be evaluated accordingly.
An element may be continued after use of the RC-mode only if no incorrect behavior related to that element caused the activation of RC-mode.}
Each activation of RC-mode is penalized. RC-mode is subject to the regulations in Section \ref{rc_mode}.

\subsection{Scoring}{
	\label{obstacle_scoring}

	\renewcommand*\footnoterule{} % No line above footnotes
	\newcommand{\topstrut}{\rule{0pt}{3.5ex}}
	\rowcolors{0}{white}{lightgray!40} % Alternate row colors

	Each team will start the event with a fixed number of base points.
	The base points are predetermined by the commission and will depend on the length of the track and the number of occuring elements.
	The total distance covered by the vehicle will have no influence on the scoring.

	\subsubsection{Timing}
	Each team will be given a \textbf{5 minute time limit} to complete \textbf{three laps} around the track.
	After the vehicle has completed all three laps, the attempt is over and no further points will be awarded.
	Timing of the event starts with the opening of the start box gate, described in section \ref{start_box}.

	\HighlightNew{After the 5 minutes have passed \textbf{or} the three laps are completed, the vehicle must be brought to an immediate and complete stop by switching to RC-mode.
	The vehicle must remain stationary until the commission signals that it may be removed from the track.}

	\subsubsection{Evaluation}
	When the vehicle passes one of the elements described in section \ref{elements_obstacle_evasion} and \ref{elements_suburban_scenario}, it will receive either a positive, neutral or negative evaluation.

	Any time the vehicle complies with all the requirements of a certain element, it will receive a positive evaluation and gain points.
	If the vehicle fails to comply with any of the requirements, it will generally receive a neutral evaluation and no points will be awarded.
	In some cases, the vehicle may receive a negative evaluation and points will be deducted.

	The following section \ref{obstacle_scoring_guidelines} provides an overview of the scoring guidelines for each element.
	The bottom row of each table shows the amount of points that will be awarded or deducted for each evaluation.

	The final score of each team will be the sum of the base points and the points 	gained or lost during the event.
	Only after a vehicle has completed at least one full lap will the attempt be valid and any points awarded.
	In case the vehicle has left the track or skipped certain parts of it, the commission will 	decide whether the attempt was valid.

	\newpage

	\subsection{Scoring Guidelines}
	\label{obstacle_scoring_guidelines}

	\subsubsection*{Static Obstacles}
	\begin{table}[H]
		\begin{tabularx}{\textwidth}{XXX}
			\toprule
			\textbf{Positive}                & \textbf{Neutral}      & \textbf{Negative}          \\
			\midrule
			Successfully passed the obstacle & Merging distance > 2m & Collision with an obstacle \\
			\topstrut
			\textbf{+10}                     & \textbf{0}            & \textbf{-10}               \\
			\bottomrule
		\end{tabularx}
	\end{table}

	\subsubsection*{Dynamic Obstacles}
	\begin{table}[H]
		\begin{tabularx}{\textwidth}{XXX}
			\toprule
			\textbf{Positive}                & \textbf{Neutral}                       & \textbf{Negative}          \\
			\midrule
			Successfully passed the obstacle & Merging distance > 2m                  & Collision with an obstacle \\
			                                 & Passed the obstacle in an intersection &                            \\
			                                 & Did not pass the obstacle              &                            \\
			\topstrut
			\textbf{+15}                     & \textbf{0}                             & \textbf{-10}               \\
			\bottomrule
		\end{tabularx}
	\end{table}

	\subsubsection*{Obstacles at Intersections}
	\begin{table}[H]
		\begin{tabularx}{\textwidth}{XXX}
			\toprule
			\textbf{Positive}                           & \textbf{Neutral}                                & \textbf{Negative}          \\
			\midrule
			Respected the right-of-way \footnotemark[2] & Did not respect right-of-way \footnotemark[2]   & Collision with an obstacle \\
			                                            & Obstacle has not fully cleared the intersection &                            \\
			\topstrut
			\textbf{+10}                                & \textbf{0}                                      & \textbf{-10}               \\
			\bottomrule
		\end{tabularx}
	\end{table}

	\subsubsection*{Oncoming Obstacles}
	\begin{table}[H]
		\begin{tabularx}{\textwidth}{XXX}
			\toprule
			\textbf{Positive}                           & \textbf{Neutral}                              & \textbf{Negative}          \\
			\midrule
			Respected the right-of-way \footnotemark[2] & Did not respect right-of-way \footnotemark[2] & Collision with an obstacle \\
			\topstrut
			\textbf{+5}                                 & \textbf{0}                                    & \textbf{-10}               \\
			\bottomrule
		\end{tabularx}
	\end{table}

	\footnotetext[2]{If the obstacle has the right-of-way, e.g. at a barred area, priority-to-right, ...}
	\newpage

	\subsubsection*{Drive Straight Stop/Give-Way Intersection}
	\begin{table}[H]
		\begin{tabularx}{\textwidth}{XXX}
			\toprule
			\textbf{Positive}                      & \textbf{Neutral}                           & \textbf{Negative}                              \\
			\midrule
			Stopped at the stop or give-way line   & Distance from stop or give-way line > 15cm & Did not stop at the stop line \footnotemark[3] \\
			Went straight through the intersection & Turned at the intersection                 &                                                \\
			                                       & Stopped for < 3s at stop line              &                                                \\
			                                       & Stopped for < 1s at give-way line          &                                                \\
			\topstrut
			\textbf{+5}                            & \textbf{0}                                 & \textbf{-5}                                    \\
			\bottomrule
		\end{tabularx}
	\end{table}

	\subsubsection*{Drive Straight Priority-to-Right Intersection}
	\begin{table}[H]
		\begin{tabularx}{\textwidth}{XXX}
			\toprule
			\textbf{Positive}                      & \textbf{Neutral}                   & \textbf{Negative}                    \\
			\midrule
			Stopped at the give-way line           & Distance from give-way line > 15cm &                                      \\
			Went straight through the intersection & Turned at the intersection         &                                      \\
			                                       & Stopped for < 1s at give-way line  &                                      \\

			\topstrut
			\textbf{+10}                           & \textbf{0}                         & \textit{No negative points possible} \\
			\bottomrule
		\end{tabularx}
	\end{table}

	\subsubsection*{Mandatory Turn Stop/Give-Way Intersection}
	\begin{table}[H]
		\begin{tabularx}{\textwidth}{XXX}
			\toprule
			\textbf{Positive}                    & \textbf{Neutral}                           & \textbf{Negative}                              \\
			\midrule
			Stopped at the stop or give-way line & Distance from stop or give-way line > 15cm & Did not stop at the stop line \footnotemark[3] \\
			Took the correct turn                & Made a wrong turn \footnotemark[4]         &                                                \\
			                                     & Stopped for < 3s at stop line              &                                                \\
			                                     & Stopped for < 1s at give-way line          &                                                \\
			\topstrut
			\textbf{+15}                         & \textbf{0}                                 & \textbf{-5}                                    \\
			\bottomrule
		\end{tabularx}
	\end{table}

	\subsubsection*{Mandatory Turn Priority-to-Right Intersection}
	\begin{table}[H]
		\begin{tabularx}{\textwidth}{XXX}
			\toprule
			\textbf{Positive}        & \textbf{Neutral}                   & \textbf{Negative}                    \\
			\midrule
			Stopped at give-way line & Distance from give-way line > 15cm &                                      \\
			Took the correct turn    & Made a wrong turn \footnotemark[4] &                                      \\
			                         & Stopped for < 1s at give-way line  &                                      \\
			\topstrut
			\textbf{+15}             & \textbf{0}                         & \textit{No negative points possible} \\
			\bottomrule
		\end{tabularx}
	\end{table}

	\footnotetext[3]{Only applies for stop and not for give-way line}
	\footnotetext[4]{Includes driving straight incorrectly}
	\newpage

	\subsubsection*{Parking}
	\begin{table}[H]
		\begin{tabularx}{\textwidth}{XXX}
			\toprule
			\textbf{Positive}                                        & \textbf{Neutral}                               & \textbf{Negative}          \\
			\midrule
			Parked correctly in the parking lot                      & Not parked correctly                           & Collision with an obstacle \\
			All wheels are within the boundaries of the parking spot & No or more than one parking maneuver performed &                            \\
			\topstrut
			\textbf{+20}                                             & \textbf{0}                                     & \textbf{-10}               \\
			\bottomrule
		\end{tabularx}
	\end{table}

	\subsubsection*{No-Passing Zones}
	\begin{table}[H]
		\begin{tabularx}{\textwidth}{XXX}
			\toprule
			\textbf{Positive}                   & \textbf{Neutral}            & \textbf{Negative}                                    \\
			\midrule
			Did not cross the solid center line & Distance to obstacle < 30cm & Crossed the solid center line with more than 1 wheel \\
			Distance to obstacle > 30cm         &                             &                                                      \\
			\topstrut
			\textbf{+5}                         & \textbf{0}                  & \textbf{-10}                                         \\
			\bottomrule
		\end{tabularx}
	\end{table}

	\subsubsection*{Barred Area}
	\begin{table}[H]
		\begin{tabularx}{\textwidth}{XXX}
			\toprule
			\textbf{Positive}             & \textbf{Neutral}        & \textbf{Negative}                    \\
			\midrule
			Did not enter the barred area & Entered the barred area &                                      \\
			\topstrut
			\textbf{+10}                  & \textbf{0}              & \textit{No negative points possible} \\
			\bottomrule
		\end{tabularx}
	\end{table}

	\subsubsection*{Crosswalk with Pedestrian(s)}
	\begin{table}[H]
		\begin{tabularx}{\textwidth}{XXX}
			\toprule
			\textbf{Positive}                      & \textbf{Neutral}                                 & \textbf{Negative}           \\
			\midrule
			Stopped at the crosswalk               & Stopped for < 3s at the crosswalk                & Collision with a pedestrian \\
			Pedestrians have cleared the crosswalk & Distance from the crosswalk > 15cm               &                             \\
			                                       & Pedestrians have not fully cleared the crosswalk &                             \\
			\topstrut
			\textbf{+15}                           & \textbf{0}                                       & \textbf{-10}                \\
			\bottomrule
		\end{tabularx}
	\end{table}

	\subsubsection*{Crosswalk without Pedestrians}
	\begin{table}[H]
		\begin{tabularx}{\textwidth}{XXX}
			\toprule
			\textbf{Positive}             & \textbf{Neutral}         & \textbf{Negative}                    \\
			\midrule
			Did not stop at the crosswalk & Stopped at the crosswalk &                                      \\
			\topstrut
			\textbf{\HighlightNew{+5}}                  & \textbf{0}               & \textit{No negative points possible} \\
			\bottomrule
		\end{tabularx}
	\end{table}

	\newpage

	\subsection*{Additional Penalties}
	\begin{table}[H]
		\begin{tabular}{@{}lccc@{}}
			\toprule
			\textbf{Violation}            & \textbf{Maximum Count} &  & \textbf{Penalty}   \\
			\midrule
			Second Attempt                & 1                      &  & -0.5 x base points \\
			Active WiFi Connection        & 1                      &  & -0.5 x base points \\
			Exceeding Speed-Limit         & $\infty$               &  & -5                 \\
			Activation of RC-Mode         & $\infty$               &  & -5                 \\
			Leaving the right lane        & 10                     &  & -2                 \\
			Collision with road sign      & $\infty$               &  & -5                 \\
			Falsely using turn indicators & 10                     &  & -2                 \\
			\bottomrule
		\end{tabular}
	\end{table}
	\clearpage
}